\documentclass[../libro.tex]{subfiles}

\begin{document}

\ifSubfilesClassLoaded{\mainmatter\chapter{行列式}\clearpage}{}

\section{缺项定位}

在正式进入本节的内容前,
我要提到一类常用的简写.
设 \(a\), \(b\), \(c\), \(d\), \(\dots\) 是若干个文字.
那么, \(a = b = c\) 表示
``\(a = b\) 且 \(b = c\)''.
同理, \(a = b = c = d\) 表示
``\(a = b\) 且 \(b = c\) 且 \(c = d\)''.
类似地, 若 \(x\), \(y\), \(z\) 是实数,
则 \(x < y < z\) 表示
``\(x < y\) 且 \(y < z\)''.
自然地, \(x < y \leq z\) 表示
``\(x < y\) 且 \(y \leq z\)''.

\vspace{2ex}

我们先看一个简单的问题.

\begin{example}
    考虑文字列~I:
    \(1\), \(2\), \(3\), \(4\), \(5\),
    \(6\), \(7\), \(8\), \(9\), \(10\).
    我们去除第~\(2\), \(6\), \(8\) 个文字
    (也就是, 去除 \(2\), \(6\), \(8\)),
    且不改变文字的前后次序,
    得文字列~II:
    \(1\), \(3\), \(4\), \(5\), \(7\), \(9\), \(10\).
    试用公式表示%
    文字列~II 的文字~\(x\) 在文字列~II 的\emph{位置}
    \(f(x)\)
    (比如, 文字 \(7\) 的位置是 \(5\)).

    此事自然不难.
    分段地, 我们可写
    \begin{align*}
        f(x) =
        \begin{cases}
            1, & x = 1;  \\
            2, & x = 3;  \\
            3, & x = 4;  \\
            4, & x = 5;  \\
            5, & x = 7;  \\
            6, & x = 9;  \\
            7, & x = 10.
        \end{cases}
    \end{align*}
    不过, 分段或许是多的.
    能否减少分段的数目?
    这是可以的.
    比如,
    我们可写
    \begin{align*}
        f(x) =
        \begin{cases}
            x,     & x < 2;     \\
            x - 1, & 2 < x < 6; \\
            x - 2, & 6 < x < 8; \\
            x - 3, & 8 < x.
        \end{cases}
    \end{align*}
    我认为这个写法是更好的,
    因为它更体现本质:
    \(x\) 前缺几项, 它的位置就减几.

    若我们想进一步地使公式简单,
    那我们可作一个新记号.
    设 \(i\), \(j\) 为二个整数.
    定义
    \begin{align*}
        \rho(i, j)
        = \begin{cases}
              0, & i < j;    \\
              1, & i \geq j.
          \end{cases}
    \end{align*}
    注意, 对不相等的二个整数 \(i\), \(j\),
    \(\rho(i, j) = 0\) 相当于 \(i < j\),
    而 \(\rho(i, j) = 1\) 相当于 \(i > j\).
    利用这个 \(\rho\)-记号, 我们可方便地写出,
    文字列~II 的文字~\(x\) 在文字列~II 的位置
    \begin{align*}
        f(x) = x - (\rho(x, 2) + \rho(x, 6) + \rho(x, 8)).
    \end{align*}
    我们记
    \(m(x) = \rho(x, 2) + \rho(x, 6) + \rho(x, 8)\).
    不难看出,
    \(x < 2\) 时, \(m(x) = 0\);
    此时, \(x\) 前不缺项.
    \(2 < x < 6\) 时, \(m(x) = 1\);
    此时, \(x\) 前缺 \(1\)~项.
    \(6 < x < 8\) 时, \(m(x) = 2\);
    此时, \(x\) 前缺 \(2\)~项.
    \(8 < x\) 时, \(m(x) = 3\);
    此时, \(x\) 前缺 \(3\)~项.
    所以, 用 \(\rho\)-记号表示的 \(f(x)\)
    跟用分段表示的 \(f(x)\) 是一样的.

    还有一件事值得提.
    因为数的加法适合结合律与交换律,
    故我们可写
    \begin{align*}
        f(x)
        = {} & x - (\rho(x, 2) + \rho(x, 6) + \rho(x, 8))  \\
        = {} & x - (\rho(x, 6) + \rho(x, 8) + \rho(x, 2))  \\
        = {} & x - (\rho(x, 8) + \rho(x, 2) + \rho(x, 6))  \\
        = {} & x - (\rho(x, 6) + \rho(x, 2) + \rho(x, 8))  \\
        = {} & x - (\rho(x, 2) + \rho(x, 8) + \rho(x, 6))  \\
        = {} & x - (\rho(x, 8) + \rho(x, 6) + \rho(x, 2)).
    \end{align*}
\end{example}

为方便, 我再介绍一次 \(\rho\)-记号.

\begin{definition}[\(\rho\)-记号]
    设 \(i\), \(j\) 为二个整数.
    定义
    \begin{align*}
        \rho(i, j)
        = \begin{cases}
              0, & i < j;    \\
              1, & i \geq j.
          \end{cases}
    \end{align*}
\end{definition}

不难看出, \(i \neq j\) 时,
\(\rho(i, j) + \rho(j, i) = 1\);
我们在后面会用到它.
(当然, \(i = j\) 时, \(\rho(i, j) + \rho(j, i) = 2\).)

\vspace{2ex}

一般地, 我们有

\begin{theorem}[缺项定位]
    设 \(n\) 为高于 \(1\) 的正整数.
    设文字列~I: \(1\), \(2\), \(\dots\), \(n\).
    设 \(k\) 是低于 \(n\) 的正整数.
    设 \(i_1\), \(i_2\), \(\dots\), \(i_k\) 是%
    不超过 \(n\) 的, 且互不相同的正整数.
    去除文字列~I 的%
    第~\(i_1\), \(i_2\), \(\dots\), \(i_k\)~个文字
    (也就是, 去除 \(i_1\), \(i_2\), \(\dots\), \(i_k\)),
    且不改变文字的前后次序,
    得文字列~II.
    那么,
    文字列~II 的文字 \(x\)
    在文字列~II 的位置
    \begin{align*}
        f(x) = x - (\rho(x, i_1) + \dots + \rho(x, i_k)).
    \end{align*}
\end{theorem}

\begin{proof}
    我们先从小到大地排
    \(i_1\), \(i_2\), \(\dots\), \(i_k\)
    为
    \(i_1'\), \(i_2'\), \(\dots\), \(i_k'\).

    为方便, 我们写
    \(i_0' = 0\),
    \(i_{k+1}' = n + 1\).
    注意, \(i_0' < 1 \leq i_1'\),
    且 \(i_{k+1}' > n \geq i_k'\).

    文字列~II 的文字恰为%
    所有不等于
    \(i_1'\), \(i_2'\), \(\dots\), \(i_k'\)
    的, 且不超过 \(n\) 的正整数.
    任取文字列~II 的一个文字 \(x\).
    那么, 存在不超过 \(k\) 的非负整数 \(v\),
    使 \(i_v' < x < i_{v+1}'\).
    于是,
    \(x\) 前缺了 \(v\)~项.
    不难验证
    \begin{align*}
        \rho(x, i_1') + \dots + \rho(x, i_k') = v.
    \end{align*}
    故 \(x\) 在文字列~II 的位置
    \begin{align*}
        f(x)
        = {} &
        x - v
        \\
        = {} &
        x - (\rho(x, i_1') + \dots + \rho(x, i_k'))
        \\
        = {} &
        x - (\rho(x, i_1) + \dots + \rho(x, i_k)).
    \end{align*}
    我们用了加法的结合律与交换律.
\end{proof}

\KunAsteriskoEnEnhavtabelo
\section{排列}
\SenAsteriskoEnEnhavtabelo

\maldevigalegajxo

为方便, 我作一个新的定义.

\begin{definition}[排列]
    设 \(a_1\), \(a_2\), \(\dots\), \(a_n\) 是%
    互不相同的 \(n\)~个文字.
    我们说, 由这 \(n\)~个文字作成的任何有序的文字列
    \(a_{i_1}\), \(a_{i_2}\), \(\dots\), \(a_{i_n}\)
    (其中, \(i_1\), \(i_2\), \(\dots\), \(i_n\)
    是互不相同的不超过 \(n\) 的 \(n\)~个正整数)
    是 \(a_1\), \(a_2\), \(\dots\), \(a_n\)
    的一个\emph{排列}.
\end{definition}

不难算出 \(n\)~个互不相同的文字的排列的数目.
从 \(a_1\), \(a_2\), \(\dots\), \(a_n\) 里选一个为%
第~\(1\)~个文字,
有 \(n\)~个选法;
从剩下的 \(n-1\)~个文字里选一个为第~\(2\)~个文字,
有 \(n-1\)~个选法;
……
从剩下的 \(2\)~个文字里选一个为第~\(n-1\)~个文字,
有 \(2\)~个选法;
从剩下的 \(1\)~个文字里选一个为第~\(n\)~个文字,
有 \(1\)~个选法.
由分步乘法计数原理, 知数目为
\begin{align*}
    n \cdot (n - 1) \cdot \dots \cdot 2 \cdot 1.
\end{align*}
不过, 在本书, 这不重要.

设文字列~I: \(1\), \(2\), \(\dots\), \(n\).
设 \(i_1\), \(i_2\), \(\dots\), \(i_n\) 是
\(1\), \(2\), \(\dots\), \(n\) 的一个排列.
我们去除文字列~I 的%
第~\(i_1\), \(i_2\), \(\dots\), \(i_{n-1}\) 个文字
(也就是, 去除 \(i_1\), \(i_2\), \(\dots\), \(i_{n-1}\)),
且不改变文字的前后次序,
得文字列~II: \(i_n\)
(恰剩一个文字).
显然, \(i_n\) 的位置 \(f(i_n) = 1\).
另一方面, 利用缺项定位公式, 有
\begin{align*}
    f(i_n) = i_n - (\rho(i_n, i_1) + \dots + \rho(i_n, i_{n-1})).
\end{align*}
于是, 我们有

\begin{theorem}
    设
    \(i_1\), \(i_2\), \(\dots\), \(i_n\)
    是
    \(1\), \(2\), \(\dots\), \(n\)
    的一个排列;
    也就是, 设
    \(i_1\), \(i_2\), \(\dots\), \(i_n\)
    是不超过 \(n\) 的正整数, 且互不相同.
    则
    \begin{align*}
        i_n - (\rho(i_n, i_1) + \rho(i_n, i_2)
        + \dots + \rho(i_n, i_{n-1})) = 1,
    \end{align*}
    或
    \begin{align*}
        \rho(i_n, i_1) + \rho(i_n, i_2)
        + \dots + \rho(i_n, i_{n-1}) = i_n - 1.
    \end{align*}
\end{theorem}

\section{\texorpdfstring{\(-1\)~的整数次方}%
  {-1 的整数次方}}

我想讲 \(-1\) 的整数次方.
在证明行列式的公式时,
我们会常用它的性质.

设 \(i\), \(j\) 为整数.
首先, 我们知道,
\begin{align*}
    (-1)^{i+j} = (-1)^{i} (-1)^{j}.
\end{align*}
因为 \((\pm 1) \cdot (\pm 1) = 1\), 故
\begin{align*}
    (-1)^{i} = (-1)^{-i}.
\end{align*}
所以
\begin{align*}
    (-1)^{i+j} = (-1)^{i} (-1)^{j}
    = (-1)^{i} (-1)^{-j}
    = (-1)^{i-j}.
\end{align*}
并且
\begin{align*}
    (-1)^{i-j} = (-1)^{-(i-j)} = (-1)^{-i+j}
    = (-1)^{-i-j}.
\end{align*}
值得提,
\begin{align*}
    (-1)^{2j} = 1.
\end{align*}
所以
\begin{align*}
    (-1)^{i \pm 2j} = (-1)^{i}.
\end{align*}

我们知道, 任给一个偶数 \(n\),
有一个整数 \(k\) 使 \(n = 2k\)
(这是定义);
任给一个奇数 (也就是, 不是偶数的整数) \(n'\),
有一个整数 \(k'\) 使 \(n' = 2k' + 1\).
所以,
若整数 \(m\) 是偶数, 则 \((-1)^m = 1\);
若整数 \(m\) 是奇数, 则 \((-1)^m = -1\).
由此可知,
若整数 \(m\) 适合 \((-1)^m = 1\),
则 \(m\) 是偶数;
若整数 \(m\) 适合 \((-1)^m = -1\),
则 \(m\) 是奇数.

最后, 我想说,
对任何的整数 \(n\),
\(n(n-1)\) 是一个偶数.
由此可见, \((-1)^n = (-1)^{n^2}\).

\KunAsteriskoEnEnhavtabelo
\section{逆序数与符号}
\SenAsteriskoEnEnhavtabelo

\maldevigalegajxo

我们知道, 排列是特别的文字列:
排列的每个文字是互不相同的.
本节, 我想介绍一些跟排列有关的量.

在研究排列时, 我们通常选文字为整数.
我们知道, 整数有大小关系.
设 \(a_1\), \(a_2\), \(\dots\), \(a_n\) 是%
互不相同的 \(n\)~个整数.
我们从小到大地排
\(a_1\), \(a_2\), \(\dots\), \(a_n\)
为
\(b_1\), \(b_2\), \(\dots\), \(b_n\)
(\(b_1 < b_2 < \dots < b_n\)).
我们说, 这是整数 \(a_1\), \(a_2\), \(\dots\), \(a_n\)
的\emph{自然排列}.
\(a_1\), \(a_2\), \(\dots\), \(a_n\)
的自然排列有且只有一个.

设
\(c_1\), \(c_2\), \(\dots\), \(c_n\)
是 \(a_1\), \(a_2\), \(\dots\), \(a_n\) 的一个排列.
若存在整数 \(i\), \(j\),
使 \(1 \leq i < j \leq n\),
且 \(c_i > c_j\),
我们说, 有序对 \((c_i, c_j)\) 为%
排列 \(c_1\), \(c_2\), \(\dots\), \(c_n\)
的一个\emph{逆序}.
比如, 排列 \(1\), \(2\), \(3\), \(4\), \(5\) 无逆序,
而 \(3\), \(2\), \(5\), \(1\), \(4\) 有 \(5\) 个逆序:
\((3, 2)\), \((3, 1)\), \((2, 1)\), \((5, 1)\), \((5, 4)\).

不难看出,
说 \((c_i, c_j)\) (其中, \(i < j\))
是排列
\(c_1\), \(c_2\), \(\dots\), \(c_n\)
的一个逆序,
相当于说 \(\rho(c_i, c_j) = 1\).
(注意, 因为
\(c_1\), \(c_2\), \(\dots\), \(c_n\)
是互不相同的整数,
故不可能出现
\(i < j\),
但 \(c_i = c_j\) 的情形.)

我们说, 一个排列的所有的逆序的数目为它的\emph{逆序数}.
比如, 排列 \(1\), \(2\), \(3\), \(4\), \(5\) 的逆序数为 \(0\),
而 \(3\), \(2\), \(5\), \(1\), \(4\) 的逆序数为 \(5\).

不难利用 \(\rho\)-记号写出排列
\(c_1\), \(c_2\), \(\dots\), \(c_n\)
的逆序数
\begin{align*}
    \tau(c_1, c_2, \dots, c_n)
    = {} & \hphantom{{} + {}} \rho(c_1, c_2)
    \\
         & + \rho(c_1, c_3) + \rho(c_2, c_3)
    \\
         & + \rho(c_1, c_4) + \rho(c_2, c_4)
    + \rho(c_3, c_4)
    \\
         & +
    \dots
    \\
         & + \rho(c_1, c_n) + \rho(c_2, c_n) + \dots
    + \rho(c_{n-1}, c_n)
    \\
    = {} & \sum_{j = 2}^{n} {
    \sum_{i = 1}^{j - 1} {\rho(c_i, c_j)} }
    \\
    = {} & \sum_{1 \leq i < j \leq n} {\rho(c_i, c_j)}.
\end{align*}
比如,
当 \(c_1\), \(c_2\), \(c_3\), \(c_4\), \(c_5\) 为
\(1\), \(2\), \(3\), \(4\), \(5\) 时,
每个 \(\rho(c_i, c_j)\)
(\(1 \leq i < j \leq 5\))
都是 \(0\),
故 \(\tau (1, 2, 3, 4, 5) = 0\);
当 \(c_1\), \(c_2\), \(c_3\), \(c_4\), \(c_5\) 为
\(3\), \(2\), \(5\), \(1\), \(4\) 时,
\(\rho(c_1, c_2)\), \(\rho(c_1, c_4)\),
\(\rho(c_2, c_4)\),
\(\rho(c_3, c_4)\), \(\rho(c_3, c_5)\)
都是 \(1\),
而
\(\rho(c_1, c_3)\), \(\rho(c_1, c_5)\),
\(\rho(c_2, c_3)\), \(\rho(c_2, c_5)\),
\(\rho(c_4, c_5)\)
都是 \(0\),
故 \(\tau (3, 2, 5, 1, 4) = 5\).

有一件小事值得提.
若一个排列恰含一个文字, 我们说, 它没有逆序,
从而它的逆序数为 \(0\).
这跟前面的公式是一样的:
既然没有整数对 \((i, j)\) 适合 \(1 \leq i < j \leq 1\),
那这个和就是 \(0\).

\vspace{2ex}

利用 \(-1\)~的整数次方,
我们可再引入一些跟排列有关的量.

\begin{definition}[符号记号]
    设 \(x\) 为整数.
    定义
    \begin{align*}
        \operatorname{sgn} {(x)}
        =
        \begin{cases}
            1,  & x > 0; \\
            0,  & x = 0; \\
            -1, & x < 0.
        \end{cases}
    \end{align*}
\end{definition}

\(\operatorname{sgn}\) 来自%
英语 \angla{sign} (或 \angla{signum}).

不难看出, 若 \(u\), \(v\) 是不相等的整数,
则 \(\operatorname{sgn} {(v - u)} = (-1)^{\rho(u, v)}\).
并且, 说 \((c_i, c_j)\) (其中, \(i < j\))
是排列
\(c_1\), \(c_2\), \(\dots\), \(c_n\)
的一个逆序,
相当于说
\(\operatorname{sgn} {(c_j - c_i)} = -1\).

\begin{definition}[整数文字列的符号]
    设 \(c_1\), \(c_2\), \(\dots\), \(c_n\) 是一个文字列,
    且它的文字全为整数
    (也就是, \(c_1\), \(c_2\), \(\dots\), \(c_n\) 是%
    整数文字列).
    定义它的\emph{符号}为
    \begin{align*}
        s(c_1, c_2, \dots, c_n)
        = {} & \hphantom{\cdot\,\,}
        \operatorname{sgn} {(c_2 - c_1)}
        \\
             & \cdot \operatorname{sgn} {(c_3 - c_1)}
        \cdot \operatorname{sgn} {(c_3 - c_2)}
        \\
             & \cdot \operatorname{sgn} {(c_4 - c_1)}
        \cdot \operatorname{sgn} {(c_4 - c_2)}
        \cdot \operatorname{sgn} {(c_4 - c_3)}
        \\
             & \cdot
        \dots
        \\
             & \cdot \operatorname{sgn} {(c_n - c_1)}
        \cdot \operatorname{sgn} {(c_n - c_2)} \cdot \dots
        \cdot \operatorname{sgn} {(c_n - c_{n-1})}
        \\
        = {} & \prod_{j = 2}^{n} {
        \prod_{i = 1}^{j - 1} {\operatorname{sgn} {(c_j - c_i)}} }
        \\
        = {} & \prod_{1 \leq i < j \leq n}
        {\operatorname{sgn} {(c_j - c_i)}}.
    \end{align*}
\end{definition}

\(s\) 也来自%
英语 \angla{sign} (或 \angla{signum}).

显然, 若 \(c_1\), \(c_2\), \(\dots\), \(c_n\) 里%
有二个相同的文字
(也就是, 存在整数 \(i\), \(j\),
使 \(1 \leq i < j \leq n\),
且 \(c_i = c_j\)),
则此文字列的符号为 \(0\).
反过来, 因为非零的整数的积不是 \(0\),
故符号为 \(0\) 的文字列%
有二个相同的文字.

有一件小事值得提.
若一个文字列恰含一个文字,
我们说, 它 (作为文字列) 的符号是 \(1\).
这跟前面的公式是一样的:
既然没有整数对 \((i, j)\) 适合 \(1 \leq i < j \leq 1\),
那这个积就是 \(1\).

不要混淆 \(\operatorname{sgn}\) 跟 \(s\):
比如, \(s(0) = 1\), 但 \(\operatorname{sgn} {(0)} = 0\);
再比如, \(s(-3) = 1\), 但 \(\operatorname{sgn} {(-3)} = -1\).

排列是特别的文字列
(排列的文字互不相同),
故排列也有符号,
其要么是 \(1\), 要么是 \(-1\).
根据 \(\operatorname{sgn}\) 与 \(\rho\) 的关系,
再利用 \(-1\) 的整数次方的性质,
不难得到排列的逆序数与符号的关系:

\begin{theorem}
    设 \(a_1\), \(a_2\), \(\dots\), \(a_n\) 是%
    互不相同的 \(n\)~个整数.
    设
    \(c_1\), \(c_2\), \(\dots\), \(c_n\)
    是 \(a_1\), \(a_2\), \(\dots\), \(a_n\) 的一个排列.
    则
    \begin{align*}
        s(c_1, c_2, \dots, c_n)
        = (-1)^{\tau(c_1, c_2, \dots, c_n)}.
    \end{align*}
\end{theorem}

一个恰含一个文字的文字列的文字也是互不相同的,
故此文字列当然是一个排列.
我们说过, 恰含一个文字的排列的逆序数为 \(0\).
因为 \((-1)^0 = 1\),
故, 约定恰含一个文字的文字列 (排列) 的符号为 \(1\),
不是没有道理的.

\end{document}

This material consists of prerequisites.

Section 1 is about...
well, what should I call it in English?
"Determining the position of the old elements in the new list"?
All right, well, whatever.
The rho notation is really useful.

Section 2 is about permutations.
This section is optional.
I do not even say "permutations" (or rather, "排列")
a lot in this book.

Section 3 is about powers of -1.
This is needed in proofs of some results about determinants.

Section 4 is about
the number of inversions of a permutation,
and the sign of a permutation.
This section is optional.

A piece of good news is that the following sections do not
depend on these four sections a lot.
